\begin{statementbox}
\textbf{这条结论解决什么问题？}


\begin{center}
\begin{tikzpicture}[>=stealth,font=\small]
\draw[gray!40,rounded corners] (0,0) rectangle (14,6.6);
\node at (7,6.1) {先看定义域成对，再看函数值相等还是相反};
\begin{scope}[shift={(3.5,2.0)},x=1.3cm,y=.85cm]
\draw[->,gray] (-1.8,0)--(1.9,0) node[right]{$x$};
\draw[->,gray] (0,-.3)--(0,3.5) node[above]{$y$};
\draw[blue,thick,domain=-1.7:1.7,samples=80] plot (\x,{\x*\x});
\draw[orange!80!black,dashed] (-1,0)--(-1,1)--(1,1)--(1,0);
\fill (-1,1) circle(1.8pt) node[left=8pt]{$(-1,1)$};
\fill (1,1) circle(1.8pt) node[right=8pt]{$(1,1)$};
\node at (0,-1.0) {$y=x^2$：关于 $y$ 轴对称};
\end{scope}
\begin{scope}[shift={(10.4,3.1)},x=1.3cm,y=.75cm]
\draw[->,gray] (-1.8,0)--(1.9,0) node[right]{$x$};
\draw[->,gray] (0,-2.2)--(0,2.5) node[above]{$y$};
\draw[teal!80!black,thick,domain=-1.3:1.3,samples=80] plot (\x,{\x*\x*\x});
\draw[orange!80!black,dashed] (-1,0)--(-1,-1)--(1,1)--(1,0);
\fill (-1,-1) circle(1.8pt) node[left=8pt]{$(-1,-1)$};
\fill (1,1) circle(1.8pt) node[right=8pt]{$(1,1)$};
\node at (0,-2.65) {$y=x^3$：关于原点对称};
\end{scope}
\end{tikzpicture}

\end{center}




同一个式子，为什么有时是偶函数，有时却不是？原因是判断奇偶性要看\textbf{整个函数}：既看允许输入哪些$x$，也看输入$x$与$-x$时的结果。

\medskip
\textbf{第一步：先看能不能成对。}

设函数$f(x)$的定义域为$D$。不论挑出哪一个允许输入的$x$，它的相反数$-x$也必须允许输入，即
\[
\boxed{\text{只要 }x\in D,\text{ 就有 }-x\in D.}
\]
这叫做\textbf{定义域关于原点对称}。如果做不到，函数\textbf{一定既不是奇函数，也不是偶函数}，到这里即可停止。

\medskip
\end{statementbox}
\clearpage
\begin{statementbox}
\textbf{第二步：再看这一对输入的结果。}
\begin{itemize}[leftmargin=*,itemsep=0.35em]
\item \textbf{偶函数：}输入$x$与$-x$，得到的数\textbf{相同}。对定义域内每个$x$都有
\[\boxed{f(-x)=f(x).}\]
图像上的点$(x,y)$与$(-x,y)$总是配对，整个图像关于$y$轴对称。
\item \textbf{奇函数：}输入$x$与$-x$，得到的数\textbf{互为相反数}。对定义域内每个$x$都有
\[\boxed{f(-x)=-f(x).}\]
图像上的点$(x,y)$与$(-x,-y)$总是配对，整个图像关于原点对称。
\end{itemize}

\textbf{计算时也可以这样写：}
\[
\begin{array}{ll}
\text{偶函数：}&f(x)-f(-x)=0,\\[2pt]
\text{奇函数：}&f(x)+f(-x)=0.
\end{array}
\]
\textbf{提醒：}以上等式必须对\textbf{所有}允许输入的$x$成立，检查一对数字并不能证明整个函数有奇偶性。


\end{statementbox}
